\documentclass[11pt]{article}
% Shared preamble for Theseus user guides.
% Compiled with pdflatex. See docs/guides/BUILD.md for rationale.

\usepackage[a4paper,margin=0.85in]{geometry}
\usepackage[T1]{fontenc}
\usepackage[utf8]{inputenc}

% Body: Palatino (via mathpazo). Headings: Helvetica (sans).
% Palatino is requested in the house style as a substitute for Charter
% because Charter is not part of the standard TeX Live distribution
% that the build container ships.
\usepackage{mathpazo}
\usepackage[scaled=0.92]{helvet}
\usepackage{courier}
\renewcommand{\familydefault}{\rmdefault}

\usepackage{microtype}
\usepackage{parskip}
\usepackage{enumitem}
\usepackage{booktabs}
\usepackage{longtable}
\usepackage{titlesec}
\usepackage{xcolor}
\usepackage{fancyhdr}
\usepackage{fancyvrb}
\usepackage{graphicx}
% Allow `_` in text mode without escaping. The screenshot file names and
% some environment-variable references in the guides contain underscores.
\usepackage{underscore}
\usepackage{tcolorbox}
\tcbuselibrary{breakable,skins}
\usepackage[hidelinks,colorlinks=true,linkcolor=black,urlcolor=teal,citecolor=black]{hyperref}

% Sans for chapter and section headings; serif body keeps prose readable
% on screen and in print.
\titleformat{\section}{\sffamily\Large\bfseries}{\thesection}{0.7em}{}
\titleformat{\subsection}{\sffamily\large\bfseries}{\thesubsection}{0.6em}{}
\titleformat{\subsubsection}{\sffamily\normalsize\bfseries}{\thesubsubsection}{0.5em}{}

\pagestyle{fancy}
\fancyhf{}
\fancyhead[L]{\small\sffamily Theseus User Guide}
\fancyhead[R]{\small\sffamily\thepage}
\renewcommand{\headrulewidth}{0.3pt}

\setlength{\parindent}{0pt}

% Convenience commands shared by every guide.
\newcommand{\page}[1]{\textbf{#1}}
\newcommand{\file}[1]{\texttt{#1}}
\newcommand{\route}[1]{\texttt{#1}}
\newcommand{\env}[1]{\texttt{#1}}
\newcommand{\code}[1]{\texttt{#1}}

% Callout boxes used by every guide.
\newtcolorbox{tldr}{
  colback=teal!5!white,colframe=teal!55!black,
  fonttitle=\sffamily\bfseries,title=TL;DR,
  breakable,boxrule=0.4pt,arc=2pt
}
\newtcolorbox{youwillneed}{
  colback=gray!4!white,colframe=gray!55!black,
  fonttitle=\sffamily\bfseries,title=You will need,
  breakable,boxrule=0.4pt,arc=2pt
}
\newtcolorbox{notcovered}{
  colback=orange!5!white,colframe=orange!55!black,
  fonttitle=\sffamily\bfseries,title=This guide does NOT cover,
  breakable,boxrule=0.4pt,arc=2pt
}
\newtcolorbox{preview}{
  colback=yellow!8!white,colframe=yellow!55!black,
  fonttitle=\sffamily\bfseries,title=Preview --- partially shipped,
  breakable,boxrule=0.4pt,arc=2pt
}
\newtcolorbox{nextup}{
  colback=blue!4!white,colframe=blue!55!black,
  fonttitle=\sffamily\bfseries,title=Where to read next,
  breakable,boxrule=0.4pt,arc=2pt
}

% Insert a screenshot only when the file exists, so the build does not
% break before the Playwright capture run is wired up.
\newcommand{\screenshot}[3]{%
  \IfFileExists{screenshots/#1}{%
    \begin{figure}[h]\centering
      \includegraphics[width=\linewidth]{screenshots/#1}
      \caption{#2}\label{fig:#3}
    \end{figure}%
  }{%
    \begin{center}\small\itshape%
      [screenshot \texttt{screenshots/#1} not yet captured --- #2]%
    \end{center}%
  }%
}


\title{\sffamily\bfseries Forecasts and Portfolio\\[2pt]\large Guide 5 of 6}
\date{}
\author{}

\begin{document}
\maketitle
\thispagestyle{empty}

\begin{tldr}
Forecasts is the prediction-market surface. The firm watches markets
on Polymarket and Kalshi, retrieves relevant sources from its own
corpus, builds a deterministic decision trace for each market, and
(only with explicit per-prediction and per-bet authorization from
founder-alpha) places a paper or live bet. Live trading is guarded by
eight successive human gates. After resolution, every prediction is
Brier-scored, log-loss-scored, and added to the public calibration
manifest. This guide explains paper vs.\ live, the Market Decision
Trace, the eight-gate safety architecture (so you understand what's
behind every public forecast), the Portfolio surface, and the
resolution / override / mismatch / revision chain. Live authorization
itself is operator-only and is covered in Guide 6.
\end{tldr}

\begin{youwillneed}
\begin{itemize}[leftmargin=*,itemsep=0.2em,topsep=0.2em]
  \item A signed-in Codex session.
  \item Familiarity with how conclusions work (Guide 2), because the
        forecast generator retrieves them as its primary inputs.
\end{itemize}
\end{youwillneed}

\begin{notcovered}
\begin{itemize}[leftmargin=*,itemsep=0.2em,topsep=0.2em]
  \item Credential storage and the per-bet operator confirmation flow
        --- operator-only, see Guide 6.
  \item Scheduler operations and the kill switch --- Guide 6.
\end{itemize}
\end{notcovered}

\section{What Forecasts and Portfolio are}

\page{Forecasts} (\route{/forecasts}) is a grid of prediction-market
opportunities the firm has scored. Each card is one prediction. A
prediction is publicly visible only after a separate authorization
step that founder-alpha performs.

\page{Portfolio} (\route{/forecasts/portfolio} and the dashboard rail)
is the live and paper position view: open positions, realized PnL,
exchange balances, exposure relative to caps.

\screenshot{05_forecasts_grid.png}{The Forecasts grid, with one open prediction in detail view.}{fc-grid}

\section{Where the markets come from}

\begin{itemize}[leftmargin=*,itemsep=0.2em]
  \item \textbf{Polymarket and Kalshi.} The system pulls markets from
        both venues on a cycle. Each market becomes a row keyed on
        venue and external id.
  \item \textbf{Operator-curated watchlist.} The firm also keeps a
        watchlist of specific market URLs to cover next cycle, which
        founder-alpha maintains.
\end{itemize}

\section{How a forecast is generated}

For each market the firm decides to score, the workshop:

\begin{enumerate}[leftmargin=*,itemsep=0.2em]
  \item Retrieves the firm's most relevant conclusions and claims via
        embedding search. The retrieval bundle must contain at least
        three distinct conclusions; otherwise the workshop abstains.
  \item Checks the near-duplicate window. If a forecast on a similar
        market was published in the last 24 hours, the workshop
        abstains.
  \item Checks the market-close buffer. If the market closes within
        one hour, the workshop refuses to predict. Stale-but-still-
        open markets are the firm's most expensive losses.
  \item Wraps the retrieval bundle in unambiguous markers and calls a
        language model with a strict JSON contract: probability,
        confidence band, a short headline, a reasoning body,
        uncertainty notes, and citations.
  \item Validates every quoted span against the cited source ---
        verbatim, character-by-character. Fabrications cause the
        prediction to abstain.
  \item Builds a \page{Market Decision Trace} (next section).
  \item Persists the prediction with status \texttt{PUBLISHED}, a
        trace carrying the decision metrics, and one citation row per
        quoted span.
\end{enumerate}

\section{The Market Decision Trace}

This is the artifact most worth understanding on the Forecasts side.
After the language model has produced a probability and a written
rationale, the workshop builds a separate object --- the \page{Market
Decision Trace} --- that is \emph{fully deterministic}: it computes a
small set of metrics from the inputs, runs a fixed rule graph over
them, and produces an action (\texttt{HOLD}, \texttt{WATCH},
\texttt{PAPER}, or \texttt{LIVE}) and a stake recommendation. No
randomness, no further model calls.

The trace records:

\begin{itemize}[leftmargin=*,itemsep=0.2em]
  \item The decision metrics it computed --- edge estimate (firm
        probability vs.\ market price), confidence, locality (how
        on-domain the retrieval was), liquidity, contradiction load,
        decay status.
  \item Each rule's firing --- which threshold was crossed, which
        veto triggered, which combination escalated to the next tier.
  \item The final action and stake.
  \item A version string (so a later refactor that changes the trace
        format is detectable and not silently mixed with old data).
\end{itemize}

The deterministic decision trace is the artifact the firm is willing
to defend. The model's prose explains, never overrides. If the prose
says ``buy'' but the trace says \texttt{WATCH}, the trace wins.

\screenshot{05_decision_trace.png}{A Market Decision Trace, with the computed metrics, the rule firings, and the resulting action.}{fc-trace}

\section{Paper vs.\ live}

The system has two trading modes:

\begin{itemize}[leftmargin=*,itemsep=0.3em]
  \item \textbf{Paper.} The default. The scheduler ingests markets,
        scores them, and writes bets to a paper-mode ledger. No order
        ever reaches an exchange.
  \item \textbf{Live.} Requires a master live-trading flag to be on
        plus the seven other gates below. Real funds move. Live mode
        is operator-only.
\end{itemize}

\section{The eight-gate safety architecture}

Whether or not you are the person clicking through them, it helps to
know the gates exist. Each gate is a separate human action;
\emph{no gate auto-promotes the next}.

\begin{enumerate}[leftmargin=*,itemsep=0.3em]
  \item \textbf{Exchange credentials configured.} Polymarket and/or
        Kalshi credentials present on the running processes.
  \item \textbf{Scheduler ingesting \& monitoring.} The forecasts
        worker is running and writing recent rows. Absence of fresh
        rows blocks every downstream step.
  \item \textbf{Paper mode validated.} Paper bets have been written,
        scored, and reconciled for long enough that the firm trusts
        the decision pipeline.
  \item \textbf{Risk caps configured.} The maximum per-bet stake and
        the maximum daily loss are set to numbers the firm is willing
        to lose. The submitter checks both on every bet.
  \item \textbf{Master live flag.} The platform-wide live-trading
        flag must be on. Without this, no other authorization counts.
  \item \textbf{Per-prediction live authorization.} Founder-alpha
        sets a flag on a specific forecast row. A row that is live-
        eligible in principle but not yet authorized is still paper-
        only.
  \item \textbf{Per-bet live confirmation.} For each individual bet,
        founder-alpha clicks confirm. The system refuses to submit
        without it.
  \item \textbf{Kill switch clear at submit time.} The submitter
        re-evaluates the kill switch immediately before placing the
        order. If the kill switch flipped between confirmation and
        submission, the bet is dropped.
\end{enumerate}

If any gate fails, every later gate refuses. The staircase is the
control surface.

\screenshot{05_forecasts_authorization.png}{The per-bet confirmation modal: live amount, gate state, and the explicit human-confirmation click target. Visible only to founder-alpha.}{fc-auth}

The reason for eight gates instead of one is to make it impossible to
accidentally place a real-money bet through software --- every gate
is an explicit human decision, and the eighth gate (kill switch
re-checked at submit time) is itself a fresh re-evaluation, not a
cached one. The integrity of the contract is also enforced inside the
schema: the database rejects any live bet that doesn't carry a
matching per-prediction authorization timestamp.

\section{The Portfolio surface}

\page{Portfolio} (\route{/forecasts/portfolio}) shows:

\begin{itemize}[leftmargin=*,itemsep=0.2em]
  \item Open positions per exchange, with side, stake, entry price,
        and current price.
  \item Realized PnL by exchange and overall.
  \item Exposure relative to the configured caps.
  \item Recent fills.
  \item Rolling 90-day Brier score and mean log-loss.
  \item The kill switch state and the daily-loss cap state.
\end{itemize}

\screenshot{05_portfolio_summary.png}{The Portfolio summary card on the founder dashboard.}{po-summary}

The portfolio surface is read-only with respect to live trading. You
cannot place a bet from the portfolio view --- bets are placed from
the per-forecast card by founder-alpha.

\begin{preview}
The portfolio's equities sub-surface (real-money equity orders) is
gated identically to the prediction-market surface but currently
ships in read-only mode for live accounts. The eight gates are wired
in; the equity submitter is staged behind an additional flag and
remains operator-only.
\end{preview}

\section{Resolution and calibration}

When a market closes, the workshop polls the venue and writes a
resolution carrying:

\begin{itemize}[leftmargin=*,itemsep=0.2em]
  \item The market outcome: \texttt{YES}, \texttt{NO},
        \texttt{CANCELLED}, or \texttt{AMBIGUOUS}.
  \item The Brier score: $(p - y)^2$, where $p$ is the firm's
        predicted probability and $y$ is 1 (YES) or 0 (NO).
  \item The log-loss: $-[y \ln p + (1-y) \ln(1-p)]$.
  \item The calibration bucket (rounded predicted probability, one
        decimal).
  \item A justification and the raw settlement payload.
\end{itemize}

If a founder believes the venue resolved incorrectly (uncommon but
real), the chain handles it:

\begin{itemize}[leftmargin=*,itemsep=0.2em]
  \item \textbf{Resolution override} --- the firm records an
        alternative resolution with a citation and reason. Carries the
        actor and timestamp.
  \item \textbf{Resolution mismatch} --- if the venue later disagrees
        with the override, a mismatch row is logged.
  \item \textbf{Resolution revision} --- an append-only revision
        history, so prior settlement is preserved for audit.
\end{itemize}

\section{The public calibration manifest}

The firm publishes its own track record. \route{/calibration} renders
a public calibration manifest --- per bucket, the number of
predictions and the realized YES rate. \route{/calibration/horizon}
shows the same broken down by time-to-resolution. The manifest is
regenerated on a schedule.

This is one of the most uncomfortable surfaces the firm offers: it is
the firm publicly grading itself on every prediction it has staked
publicly. If the firm has been consistently overconfident at the
\texttt{0.7} bucket, the public page shows it.

\section{The kill switch and daily loss cap}

These are operator surfaces, but as a founder you should know they
exist:

\begin{itemize}[leftmargin=*,itemsep=0.2em]
  \item \textbf{Per-submit kill switch.} Re-checked immediately before
        each order placement. Trips even bets that already passed
        operator confirmation. Open positions are not auto-closed;
        they are closed manually.
  \item \textbf{Daily-loss cap.} Once tripped, the submitter refuses
        new orders for the rest of the calendar day, regardless of
        how green the other gates are. Resets by calendar day.
\end{itemize}

\section{What you actually do with Forecasts day to day}

For most founders, the Forecasts surface is a reading view:

\begin{itemize}[leftmargin=*,itemsep=0.2em]
  \item Browse the grid to see what the firm is predicting.
  \item Click a card to read its decision trace and citation chain.
        Each citation links back to a firm conclusion, so you can
        verify the chain of reasoning yourself.
  \item Watch resolutions over time; the rolling Brier score on the
        portfolio card tells you whether the firm's calibration is
        holding up.
  \item Suggest watch-list additions to founder-alpha for markets
        you'd like the firm to score.
\end{itemize}

\section{Where the docs are}

\begin{itemize}[leftmargin=*,itemsep=0.2em]
  \item The other five guides in this series.
  \item Guide 6 covers the operator surfaces for Forecasts: per-bet
        confirmation, kill switch, resolution override, and the
        equities preview.
\end{itemize}

\begin{nextup}
Continue with \href{06_Operator_Console.pdf}{Guide 6 --- Operator
Console} for what founder-alpha does behind the scenes: kill
switches, live authorization, the publication review queue, the
quarterly methodology review week, deletion vs.\ retraction. Reading
Guide 6 is useful even if you don't have operator privileges --- it
explains the gates that everything else passes through.
\end{nextup}

\end{document}
